\documentclass{article}
\input{preamble.pre}
\begin{document}

% Negative dimensions directly in a vertical list: unlike inside an hbox row or a paragraph
% line (where other material keeps the line's own height/depth >= 0), a negative total here
% pulls the following material up.
\block{Negative in vertical list}{
    \lorem\par
    \showvwith{}{-20pt}{}{\hbox{Content}}
    \lorem\par
    \showvwith{}{}{-20pt}{\hbox{Content}}
    \lorem\par
    \showvto{-20pt}{\hbox{Content}}
    \lorem\par
    \showvwith{}{-20pt}{5pt}{\hbox{Content}}
    \lorem\par
}

\block{Negative vtop/hbox in vertical list}{
    \lorem\par
    \showtwith{}{-20pt}{}{\hbox{Content}}
    \lorem\par
    \showtwith{}{}{-20pt}{\hbox{Content}}
    \lorem\par
    \showhwith{}{-10pt}{}{Content}
    \lorem\par
    \showhwith{}{}{-10pt}{Content}
    \lorem\par
}

\pagebreak

% The last row of a vbox determines its natural depth. With \ht0= or "to" the CSS has to
% approximate that depth (it uses one text line's descent), so these vary what the last item
% is: a rule (depth 0), a box with a large depth, trailing glue (depth 0). The blue brackets
% sit on the enclosing hbox's baseline.
\block{Last item: rule}{
    \showh{\boxwith{}{50pt}{}{\vbox{\hbox{Content}\kern3pt\hrule height2pt}}}
    \showh{\setbox0\vbox to 50pt{\hbox{Content}\kern3pt\hrule height2pt}\box0}
    \showh{\setbox0\vbox to 50pt{\hbox{Content}\kern3pt\hrule height2pt}\dp0=10pt\box0}
}
\block{Last item: deep box}{
    \showh{\boxwith{}{50pt}{}{\vbox{\hbox{Content}\hbox{\vrule height5pt depth20pt width20pt}}}}
    \showh{\setbox0\vbox to 50pt{\hbox{Content}\hbox{\vrule height5pt depth20pt width20pt}}\box0}
    \showh{\setbox0\vbox to 50pt{\hbox{Content}\hbox{\vrule height5pt depth20pt width20pt}}\dp0=10pt\box0}
}
\block{Last item: glue}{
    \showh{\boxwith{}{50pt}{}{\vbox{\hbox{Content}\vskip5pt}}}
    \showh{\setbox0\vbox to 50pt{\hbox{Content}\vskip5pt}\box0}
    \showh{\setbox0\vbox to 50pt{\hbox{Content}\vskip5pt}\dp0=10pt\box0}
}

\block{VBox to with negative depth}{
    \showh{\setbox0\vbox to 50pt{\showvinner{\hbox{Content}}}\dp0=-10pt\box0}
    \noindent\showpar{\setbox0\vbox to 50pt{\showvinner{\hbox{Content}}}\dp0=-10pt\box0}
}

\pagebreak

% A lone \ht0= or \dp0= SMALLER than the content: the box must claim only the asserted amount
% on that side, so neighbouring lines move closer and the content overlaps them.
\block{HBox height/depth smaller than content}{
    \showh{\showhwith{}{2pt}{}{\vrule height20pt depth5pt width5pt Content}}
    \showh{\showhwith{}{}{1pt}{\vrule height5pt depth20pt width5pt Content}}
    \noindent\showpar{\showhwith{}{2pt}{}{\vrule height20pt depth5pt width5pt Content}}
    \noindent\showpar{\showhwith{}{}{1pt}{\vrule height5pt depth20pt width5pt Content}}
}

% Asserted boxes nested in asserted boxes, inside a paragraph.
\block{Nested}{
    \noindent\showpar{\showhwith{}{5pt}{2pt}{\showvwith{}{40pt}{}{\hbox{Content}}}}
    \noindent\showpar{\showhwith{}{30pt}{}{\showtwith{}{}{40pt}{\hbox{Content}}}}
    \showh{\showvwith{}{40pt}{10pt}{\showh{\showvto{30pt}{\hbox{Content}}}}}
}

\end{document}
